\chapter{System Definition Document}

\section{Components}
\begin{enumerate}
  \item \textbf{Study pipeline:} selects assets, builds the historical dataset and runs the backtest
        outside the web application.
  \item \textbf{Data routine:} collects data outside the server on business days and generates a
        snapshot with the model inputs.
  \item \textbf{Backend:} Python API deployed on PythonAnywhere, which reads the snapshot and
        performs the model calculations.
  \item \textbf{Frontend:} separate user interface that consumes the API. Streamlit is the initial
        option and may change.
\end{enumerate}

\section{Inputs}
\begin{enumerate}
  \item \textbf{Asset universe:} assets selected by liquidity in the study, or selected by the user
        from the catalog in the application.
  \item \textbf{Market equilibrium:} market capitalization weights, used to obtain implied returns
        by reverse optimization, consistent with the CAPM.
  \item \textbf{Investor views:} absolute or relative views on asset returns, available for current
        portfolios.
\end{enumerate}

\section{Processing}
\begin{enumerate}
  \item Selection of the assets by liquidity.
  \item Collection of prices, CDI and market capitalization of the assets.
  \item Data treatment: alignment, missing values, returns and covariance.
  \item Construction of the equilibrium prior.
  \item Combination of the prior with the views (posterior).
  \item Computation of the optimal weights.
  \item Walk-forward backtest in the study, without views in the main version.
\end{enumerate}

\section{Output}
\begin{itemize}
  \item Weights table: market portfolio (prior) $\times$ BL portfolio $\times$ difference.
  \item Expected portfolio metrics: return, volatility and Sharpe ratio.
  \item Study backtest results: chart and metrics compared to the benchmarks.
\end{itemize}

\section{User requirements}
\begin{itemize}
  \item Basic knowledge of financial markets and stocks.
  \item Ability to express views as expected returns in \% per year above the CDI.
\end{itemize}
